\section*{Введение}
\addcontentsline{toc}{section}{Введение}

\textbf{Актуальность темы.}  
В современных условиях цифровизации повседневной жизни возрастает интерес к персональным веб-приложениям, направленным на упрощение управления бытовыми задачами. Одной из таких задач является организация личного гардероба. Большое количество одежды и отсутствие систематизированного подхода к её учёту приводят к сложностям при выборе повседневных образов, дублированию вещей и неэффективному использованию гардероба.

Разработка веб-приложений для управления личным гардеробом позволяет структурировать информацию об одежде, упростить её каталогизацию и автоматизировать процесс формирования комплектов одежды (образов), что повышает удобство использования и улучшает пользовательский опыт.

\textbf{Степень разработанности проблемы.}  
В настоящее время существует ряд мобильных и веб-приложений, предназначенных для управления гардеробом и подбора одежды, таких как Stylebook, Smart Closet и Cladwell. Данные решения предоставляют базовые функции каталогизации и формирования образов, однако часто ограничены в гибкости настройки и возможностях расширения логики под индивидуальные потребности пользователя.

Кроме того, большинство существующих систем ориентированы на закрытые платформы и не предоставляют открытой архитектуры для модификации или интеграции дополнительных алгоритмов.

\textbf{Объект исследования} — процессы организации и использования личного гардероба в цифровой среде.

\textbf{Предмет исследования} — методы и программные средства для автоматизации управления элементами гардероба и формирования пользовательских образов.

\textbf{Целью данной работы} является разработка веб-приложения для управления личным гардеробом и формирования пользовательских образов.
